% Generic solution template for a Sussex assessed problem set.
% Use this ONLY when the module supplies no template of its own.
% If the instructor supplies one, use theirs and leave its preamble untouched.

\documentclass[12pt]{article}

\usepackage[a4paper,margin=24mm]{geometry}
\usepackage{amsmath,amssymb,mathtools}
\usepackage{braket}
\usepackage[most]{tcolorbox}
\usepackage{enumitem}
\usepackage{graphicx}
\usepackage{xcolor}
\usepackage[T1]{fontenc}
\usepackage{lmodern}
\usepackage{microtype}
\usepackage{float}
\usepackage[hidelinks]{hyperref}

\renewcommand{\d}{\mathop{}\!\mathrm{d}}
\setlength{\parindent}{0pt}
\setlength{\parskip}{0.55em}
\graphicspath{{figures/}}
\setcounter{secnumdepth}{0}

\tcbset{enhanced, breakable, boxrule=0.8pt, arc=1.5mm,
        left=2mm, right=2mm, top=1.5mm, bottom=1.5mm}

\newtcolorbox{problembox}[1]{colback=gray!8,  colframe=black,
                             fonttitle=\bfseries, title={#1}}
\newtcolorbox{solutionbox}{  colback=blue!3,  colframe=blue!60!black,
                             fonttitle=\bfseries, title={Solution}}
\newtcolorbox{answerbox}{    colback=yellow!10, colframe=orange!70!black,
                             fonttitle=\bfseries, title={Final answer}}

% =====================================================================
%  REQUIRED SLOT: every solution ends with a closing sentence.
%
%  Tutor feedback on a 98/100 script, repeated across four questions:
%    "I appreciate the neatly boxed final answer, but it will be
%     appropriate to write a sentence of text to help close off the
%     final answer."
%
%  \closing{...}  renders that sentence.
%  \closingmissing is the placeholder. Leave it in by accident and it
%  prints a red warning into the PDF, so a page-read catches it.
% =====================================================================
\newcommand{\closing}[1]{\par\medskip\noindent\textbf{Therefore,} #1}
\newcommand{\closingmissing}{%
  \par\medskip\noindent
  \fcolorbox{red}{red!12}{\parbox{0.95\linewidth}{\bfseries
  CLOSING SENTENCE MISSING --- replace \texttt{\textbackslash closingmissing}
  with \texttt{\textbackslash closing\{...\}}: a sentence naming the result.}}}

\begin{document}

\begin{center}
{\Large\bfseries MODULE CODE -- Module Title\\[0.2em]
Assessed problem set -- Week N\\[0.3em]
{\normalsize Weight: NN\% of module grade}\\[1em]
Candidate number: REPLACE-ME}
\end{center}

% =====================================================================
\section{Question 1 (Canvas question 1)}

\begin{problembox}{Question 1 \hfill [N marks]}
Paste the question statement here, from the instructor-authored source,
not from a platform markdown export. See references/checks.md.
\end{problembox}

\begin{solutionbox}
\textbf{What is asked.} Restate the question and list the givens.

\textbf{Principle used.} Name the principle in words and say why it applies
here. Marks are explicitly allocated to this when notes are permitted.

\textbf{Working.} Show every step. Scale the prose to the marks: a 2-mark
part can be three lines of algebra, a 15-mark part needs explanation
between the equations.

\textbf{Check.} An independent check -- dimensions, a limiting case, a
probability summing to one, a sign agreeing with a sketch.
\end{solutionbox}

\begin{answerbox}
State the result.
\closingmissing   % <-- replace with \closing{...}
\end{answerbox}

\end{document}
